% from ega4-16-fr.tex
\oldpage[IV-4]{34}
Revenant au problème considéré en (16.5.13), on obtient donc:

\begin{proposition}[16.5.17]
\label{IV.16.5.17}
Soient $X$, $Y$ deux $S$-préschémas, $Y_0$ un sous-préschéma fermé de $Y$
défini par un Idéal quasi-cohérent $\mathscr{J}$ de $\mathcal{O}_Y$ tel que
$\mathscr{J}^2=0$, $j:Y_0\to Y$ l'injection canonique. Soit
$u_0:Y_0\to X$ un $S$-morphisme, et $\mathscr{P}$ le faisceau d'ensembles sur
$Y$ tel que, pour tout ouvert $U$ de $Y$, $\Gamma(U,\mathscr{P})$ soit
l'ensemble des $S$-morphismes $u:U\to X$ tels que
$u_0|U_0=u\circ(j|U_0)$, où $U_0=j^{-1}(U)$. Alors il existe sur
$\mathscr{P}$ une structure de pseudo-torseur sous le
$\mathcal{O}_{Y_0}$-Module
\[
  \mathscr{G}=\mathscr{H}om_{\mathcal{O}_{Y_0}}
  (u_0^*(\Omega_{X/S}^1),\mathscr{J}).
\]
\end{proposition}

En particulier:

\begin{corollary}[16.5.18]
\label{IV.16.5.18}
Avec les notations de (16.5.16), supposons $Y$ affine et $\Omega_{X/S}^1$ de
présentation finie; s'il existe un recouvrement ouvert $(U_\alpha)$ de $Y$,
et, pour chaque indice $\alpha$, un $S$-morphisme $v_\alpha:U_\alpha\to X$
tel que, si $U_\alpha^0=j^{-1}(U_\alpha)$, on ait
$v_\alpha\circ(j|U_\alpha^0)=u_0|U_\alpha^0$, alors il existe un
$S$-morphisme $u:Y\to X$ tel que $u\circ j=u_0$.
\end{corollary}

En effet, $\mathscr{G}$ est alors un $\mathcal{O}_{Y_0}$-Module
quasi-cohérent (I, 1.3.12); en vertu de (16.5.16) et du fait que $Y_0$ est
alors affine, le faisceau $\mathscr{P}$, qui est par hypothèse un torseur sous
$\mathscr{G}$, et non seulement un pseudo-torseur, est \emph{trivial}; mais si
$w$ est un isomorphisme de $\mathscr{G}$ sur $\mathscr{P}$ (en tant que
torseurs sous $\mathscr{G}$), l'image par $w$ de la section 0 de
$\mathscr{G}$ est le $S$-morphisme $u$ cherché.

\subsection{\texorpdfstring{Faisceaux de $p$-différentielles et différentielle extérieure.}{Faisceaux de p-différentielles et différentielle extérieure.}}
\EGAReaderTarget{ega4.c04.subsec.16-6}{IV.16.6}

\begin{env}[16.6.1]
\EGAReaderTarget{ega4.c04.statement.16-6-1}{IV.16.6.1}
Soit $f:X\to S$ un morphisme de préschémas. On appelle \emph{faisceau des
$p$-différentielles de $X$ relativement à $S$} ($p$ entier) la
\emph{puissance extérieure $p$-ème} (0, 4.1.5) du $\mathcal{O}_X$-Module
$\Omega_{X/S}^1$, notée
\[
  \Omega_{X/S}^p=\bigwedge^p(\Omega_{X/S}^1).
  \tag{16.6.1.1}
  \label{IV.16.6.1.1}
\]
On a donc $\Omega_{X/S}^0=\mathcal{O}_X$, et $\Omega_{X/S}^p=0$ pour $p<0$;
les $\Omega_{X/S}^p$ sont les composants homogènes de l'Algèbre extérieure
sur $\Omega_{X/S}^1$
\[
  \Omega_{X/S}^{\bullet}=\bigwedge(\Omega_{X/S}^1)
  =\bigoplus_{p\in\mathbf{Z}}\bigwedge^p(\Omega_{X/S}^1),
  \tag{16.6.1.2}
  \label{IV.16.6.1.2}
\]
qui est donc une $\mathcal{O}_X$-Algèbre graduée quasi-cohérente,
anticommutative et dont les éléments de degré 1 sont de carré nul. Pour tout
ouvert affine $U$ de $X$, on a
\[
  \Gamma(U,\Omega_{X/S}^{\bullet})
  =\bigwedge(\Gamma(U,\Omega_{X/S}^1)),
\]
où $\Gamma(U,\Omega_{X/S}^1)$ est considéré comme
$\Gamma(U,\mathcal{O}_X)$-module.

Lorsque $S=\operatorname{Spec}(A)$ et $X=\operatorname{Spec}(B)$ sont
affines, $B$ étant donc une $A$-algèbre, on a (0, 4.1.5)
$\Omega_{X/S}^p=(\Omega_{B/A}^p)^{\sim}$, en posant
$\Omega_{B/A}^p=\bigwedge^p\Omega_{B/A}^1$.
\end{env}

\begin{theorem}[16.6.2]
\label{IV.16.6.2}
Il existe un endomorphisme et un seul $d$ du faisceau de groupes additifs
$\Omega_{X/S}^{\bullet}$, ayant les propriétés suivantes:
\begin{enumerate}
\item[(i)] $d\circ d=0$.
\item[(ii)] Pour tout ouvert $U$ de $X$ et toute section
$f\in\Gamma(U,\mathcal{O}_X)$, on a $df=d_{X/S}f$.
